\begin{afterword}
\summaryhead

The author recalls telling a man who doubted evolution that humans and chimpanzees share 98\% of their DNA. The man suggested the similarity might be coincidence. Told that the odds against that were two to the power of 750 million to one, he replied, ``But there's still a chance, right?''

The post concedes that the figure's source is forgotten, and that odds so extreme ignore the chance that the calculation itself is wrong; the 98\% was later revised to 95\%. But it finds the reply revealing. First, intuition treats ``no chance'' and ``a very tiny chance'' as different in kind, because no feeling is small enough to represent a tiny probability. That, the post says, is why people buy lottery tickets. Second, people treat an argument that is not certain as one they may ignore. The author answers that if a certain disproof cannot be ignored, neither can a likelihood of one in a googol, or one of 0.9. The lesson applies whenever someone says, ``But you can't prove me wrong.''

\responsehead

The post's closing argument is sound. If a proof may not be ignored because one dislikes it, neither may a strong probabilistic argument, and the post applies this to the reader's own reasoning, not only to the man's. It is also candid about its anecdote: it says the author cannot source the figure and that people who state such odds are wrong far more often than the odds allow.

The first insight is better supported than the post shows. Something close to it, that people treat a tiny chance and no chance as different in kind, was part of Kahneman and Tversky's prospect theory in 1979, where decision weights change abruptly near zero and help explain gambling. The post's only evidence for it is the anecdote and the blog's earlier posts on the lottery. Its own explanation, that a feeling that small ``doesn't fire enough neurons,'' is a claim about the brain with no source. It repeats the observation in terms of neurons and adds no evidence for it.

The second insight has even less support. The rule that ``if an argument is not certain, you're allowed to ignore it'' is the author's reading of one sentence, spoken years earlier and recalled from memory. The man is not reported saying anything more. The post offers no other case.

One factual doubt the post raises about itself can be settled. It worries that the revised 95\% figure may apply only to genes, in which case its odds would be off by a ``meta-order of magnitude.'' The 95\% estimate came from comparing genomic DNA, repeated and non-repeated alike, and the 2005 comparison of the whole genomes found the same kinds of difference across the genome. So the figure is not limited to genes, and the odds are not off for that reason.

In short: a sound argument against ignoring probabilistic evidence, with a first insight that prior research supports and a second drawn from a single remembered sentence.
\end{afterword}
